\section{Coefficient five and the full small quadratic family}
\label{sec:coefficient-five-small-quadratic}

\paragraph{Original coefficient and finite character convention.}
Put $K=\mathbb Q(\omega)$ and $\eta_K=\chi_{-3}$.
Let $E/\mathbb Q_5$ be any finite extension, with its original
ring of integers $O$, containing the values of a generator
$\xi_7$ of the six Dirichlet characters modulo seven.
Choose its indexing by $\xi_7(3)=\zeta_6$, a primitive sixth
root in this specified coefficient. Replacing the generator
by its inverse does not alter the families below.
Put $D_O=E/O$, let $c_5$ be the ordinary cyclotomic character,
and let $\tau_5$ be its Teichm\"uller lift of order four.
The values of $\tau_5$ already lie in $\mathbb Z_5^\times$;
no extension of $O$ is made. Define
\[
 \begin{gathered}
 \mathcal D=\{1,\chi_{-4},\chi_8,\chi_{-8}\},\qquad N_T=T-1,\\
 \mathcal A_e=\begin{cases}
       \mathcal D\times\{0,1,\ldots,5\},&4\nmid e,\\
       \{(1,j):j\in\{0,2,4\}\}\ \sqcup\
       \{(\chi_{-4},j),(\chi_8,j):0\leq j<6\},&4\mid e.
                  \end{cases}
 \end{gathered}\tag{F5.1}
\]
For $(\delta,j)\in\mathcal A_e$ let
$\epsilon(\delta)\in\{0,1\}$ satisfy
$\delta(-1)=(-1)^{\epsilon(\delta)}$, and choose
\[
 \begin{gathered}
 b\in\{0,1\},\qquad b\equiv\epsilon(\delta)+j-e\pmod2,\\
 \psi^{\rm nc}=\delta\xi_7^j\eta_K^b,\quad
 \psi^{\rm cr}=\psi^{\rm nc}\eta_K,\qquad
 a^{\rm nc}=\psi^{\rm nc}c_5^{-e},\quad
 a^{\rm cr}=\psi^{\rm cr}c_5^{-e}.
 \end{gathered}\tag{F5.2}
\]
The finite parities are $(-1)^e$ and $(-1)^{e+1}$.
Thus these are exactly the noncritical and critical rational
extensions of (CP.3). The auxiliary $\tau_5$ used below belongs
only to the residual calculation, not to the finite host label.

For a number field $M$ write $\mathcal S_5(M,A)$ for the full
strict kernel: continuous global $H^1$ restricted to full
$H^1(M_v,A)$ at every $v\mid5$, and to $H^1(I_v,A)$ at all
other finite places. Include real-place $H^1$ if present;
it is zero on a $5$-primary coefficient. These are the
original full local conditions, including their finite
localization interpretation in (FC.2).

\begin{theorem}[AC-G199: coefficient-five small quadratic descent]
\label{AC-G199}
For every such original $O$, every integer $e\geq1$ and every
$(\delta,j)\in\mathcal A_e$,
\[
                  \mathcal S_5(\mathbb Q,D_O(a^{\rm nc}))=0.
 \tag{F5.3}
\]
On the same original coefficient,
\[
 \mathcal S_5\bigl(K,D_O(\delta\xi_7^j c_5^{-e})\bigr)
                \simeq\mathcal S_5(\mathbb Q,D_O(a^{\rm cr})).
 \tag{F5.4}
\]
The right side is not asserted zero.

For the actual host and coefficients of (FC.1)--(FC.2), with
$h\geq2$ and integer $T\geq1$, put
\[
 Y^{\rm small}_5(T)=
 \{(5,(\delta\xi_7^j)|_{G_K},e):1\leq e<T,
                                      (\delta,j)\in\mathcal A_e\}.
 \tag{F5.5}
\]
All undisplayed cubic and auxiliary factors are trivial.
These are distinct actual labels, with exact count
\[
          |Y^{\rm small}_5(T)|=24N_T-9\left\lfloor N_T/4\right\rfloor.
 \tag{F5.6}
\]
For $T\geq2$ and every subset $Y'\subseteq Y^{\rm small}_5(T)$,
\[
 \sum_{(5,\eta,e)\in Y'}
             \frac{\log\#\mathcal S_{5,\eta}(e)}{b_5}
       \leq\mathcal B_{\rm cr}(F,T)\leq C\mathcal P(F,T),
 \tag{F5.7}
\]
with exactly the one original-host aggregate of AC-G191.
For every retained set $J$ of actual triples in this weight range,
\[
                \mathcal R_J=
                      \mathcal R_{J\setminus Y^{\rm small}_5(T)}.
 \tag{F5.8}
\]
This is zero-term deletion from the noncritical rational remainder.
It does not split positive derived indices or assert original-$K$
strict zero. Old-union intersections and their window count are
separate from the internally distinct count (F5.6).
\end{theorem}

\begin{proof}
\proofstep{1. The exact even residual character.}
Put $t=e\bmod4$ and $a=j\bmod2$, represented by
$0\leq t<4$ and $a\in\{0,1\}$. The finite character
\[
                 \rho=\delta\xi_7^j\eta_K^b\tau_5^{-t}
 \tag{F5.9}
\]
has order dividing twelve and is even. Its reduction is exactly
$\overline{a^{\rm nc}}$, since $\overline{\tau_5}=\overline{c_5}$.
Reduction is injective on roots of unity of order prime to five.
Thus its kernel is the literal residual kernel in the original
residue field, with no choice of a new embedding.

Let $R=\mathbb Q(\zeta_7)^+$ and put
$P_7(Y)=Y^3+Y^2-2Y-1$.
The identity $\Phi_7(Z)=Z^3P_7(Z+Z^{-1})$ identifies its
real trace generator and degree three. The square of $\xi_7$
factors through $R$, and $\xi_7^3=\chi_{-7}$.
Strip off the even cubic power as
\[
 \rho=\xi_7^{j-3a}\theta,\qquad
              \theta=\delta\chi_{-7}^{a}\eta_K^b\tau_5^{-t}.
 \tag{F5.10}
\]
We construct explicit real covers of the field of $\theta$;
their composita with $R$ therefore kill $\rho$.

\proofstep{2. The quartic covers at both odd residues.}
For odd $t$, the quadratic factor in (F5.10) is the character
of $\mathbb Q(\sqrt{-A})$, with $A$ specified by this table:
\[
 \begin{array}{c|cc}
 \delta&A\ (a=0)&A\ (a=1)\\\hline
 1&3&7\\
 \chi_{-4}&1&21\\
 \chi_8&6&14\\
 \chi_{-8}&2&42
 \end{array}\qquad
                \theta=\chi_{\mathbb Q(\sqrt{-A})}\tau_5^{-t}.
 \tag{F5.11}
\]
The notation on the right uses the quadratic field character,
so, for example, $A=1,2,6$ have primitive discriminants
$-4,-8,-24$. The entries follow by multiplying quadratic
squareclasses in $\delta\chi_{-7}^a\eta_K^b$;
$b\equiv\epsilon(\delta)+a-1\pmod2$ here.

For a primitive fifth root $z$, put
\[
             \beta_A=\sqrt{-A}(z-z^{-1}),\qquad
                  Q_A(X)=X^4-5AX^2+5A^2.
 \tag{F5.12}
\]
The product is real, and
$\beta_A^2=A(5+\sqrt5)/2$, after choosing the sign of $\sqrt5$.
Its conjugate square is $A(5-\sqrt5)/2$, so $Q_A(\beta_A)=0$.
The fields $\mathbb Q(z)$ and $\mathbb Q(\sqrt{-A})$ are
disjoint: the only quadratic subfield of the cyclic quartic
$\mathbb Q(z)$ is the real field $\mathbb Q(\sqrt5)$.
In their group $C_4\times C_2$, the stabilizer of $\beta_A$
is exactly the subgroup generated by the simultaneous sign
changes $(z\mapsto z^{-1},\sqrt{-A}\mapsto-\sqrt{-A})$.
Indeed any stabilizer fixes $\beta_A^2$, hence fixes $\sqrt5$;
the two indicated possibilities then follow.
This subgroup has order two and is precisely the kernel of
the order-four character in (F5.11), for either odd $t$.
Consequently
\[
          F_A=\mathbb Q(\beta_A)
             \quad\hbox{is the exact real cyclic quartic field of }\theta.
 \tag{F5.13}
\]
Replacing the quartic character by its inverse leaves this field
unchanged. This verifies the field construction independently
of a numerical root or a proposed polynomial label.

\proofstep{3. The even-residue covers and all twenty exact polynomials.}
At $t=2$, $\tau_5^{-t}=\chi_5$, so $\theta$ is the real
quadratic field character of radicand $D$ in the first table.
At $t=0$ the admitted cases give the second table:
\[
 \begin{array}{c|cc}
 t=2& D\ (a=0)&D\ (a=1)\\\hline
 1&5&105\\
 \chi_{-4}&15&35\\
 \chi_8&10&210\\
 \chi_{-8}&30&70
 \end{array}
 \qquad
 \begin{array}{c|cc}
 t=0& D\ (a=0)&D\ (a=1)\\\hline
 1&\text{trivial}&\text{not admitted}\\
 \chi_{-4}&3&7\\
 \chi_8&2&42\\
 \chi_{-8}&\text{not admitted}&\text{not admitted}
 \end{array}\tag{F5.14}
\]
For instance $\chi_{-4}\eta_K\chi_5$ has radicand $15$,
and $\chi_8\chi_{-7}\eta_K\chi_5$ has radicand $210$.
Every entry follows from the same formula (F5.10) with the
stated parity $b$; the exact finite-character checker verifies
each identity at all automorphisms of $\mathbb Q(\zeta_{840})$.

The required covers $M$ are $RF_A$ for odd $t$,
$R\mathbb Q(\sqrt D)$ for the nontrivial even entries, and
$R$ for the trivial entry. The complete parameter lists are
\[
 \begin{split}
 \mathcal A&=\{1,2,3,6,7,14,21,42\},\\
 \mathcal D_0&=\{2,3,7,42,5,10,15,30,35,70,105,210\}.
 \end{split}\tag{F5.15}
\]
Their twenty defining polynomials are given exactly, without
rounding, by
\[
 \begin{aligned}
 P_A^{(4)}(X)&=\operatorname{Res}_Y(P_7(Y),Q_A(X-Y)),
                                      &&A\in\mathcal A,\\
 P_D^{(2)}(X)&=\operatorname{Res}_Y(P_7(Y),(X-Y)^2-D),
                                      &&D\in\mathcal D_0.
 \end{aligned}\tag{F5.16}
\]
They have degrees twelve and six. To verify that these are
the asserted fields, $R$ intersects either quartic or quadratic
factor in $\mathbb Q$, by coprime degrees. Let $\alpha$ be
its trace generator and let $\beta$ generate the other factor.
An automorphism fixing $\alpha+\beta$ has
$\sigma\alpha-\alpha=\beta-\sigma\beta$ in this intersection.
It is a rational number $q$. Iterating the finite-order
automorphism gives $nq=0$, hence $q=0$. Both generators are
then fixed. Thus the sum is primitive in the compositum,
and (F5.16) is its minimal polynomial.

All twenty coefficient arrays are retained in the exact input
JSON. As an independent polynomial check, multiplication by
$\alpha$ in $\mathbb Q[Y]/P_7$ has matrix
\[
                     C=\begin{pmatrix}0&0&1\\1&0&2\\0&1&-1\end{pmatrix}.
 \tag{F5.17}
\]
The determinants $\det Q_A(XI-C)$ and
$\det((XI-C)^2-DI)$ reproduce every coefficient of (F5.16).
The finite cover check also verifies the exact containment
$\ker(\text{cover})\subseteq\ker\rho$ in all admitted cases.

\proofstep{4. Successful class-quotient certificates and the cubic case.}
For each polynomial in (F5.16), a fresh \texttt{bnfinit(P,1)}
followed by \texttt{bnfcertify(B,1)} returned one.
The computed quotient groups, in increasing $A$ order, are
\[
 \begin{array}{c|rrrrrrrr}
 A&1&2&3&6&7&14&21&42\\\hline
 \text{quotient}&1&1&1&C_2&1&C_2&C_2^2&C_2^2
 \end{array}\tag{F5.18}
\]
For the quadratic composita they are
\[
 \begin{array}{c|rrrrrr}
 D&2&3&7&42&5&10\\\hline
 \text{quotient}&1&1&1&C_2&1&C_2
 \end{array}
 \quad
 \begin{array}{c|rrrrrr}
 D&15&30&35&70&105&210\\\hline
 \text{quotient}&C_2&C_2&C_2&C_2&C_2&C_2^2
 \end{array}\tag{F5.19}
\]
The applied PARI flag-one contract certifies that each actual
class group is a quotient of the displayed group. It does not
certify a complete unit basis or equality with a nontrivial
computed quotient. In particular every one of these twenty
actual fields has trivial $5$-primary class group.
The original route, cover formulas and successful class
executions were supplied by root. The exact matrices,
character identities and fresh owned replay are separately
identified in the computation receipts; they are not an
independent proof of the underlying PARI algorithm.

For $R$ itself, $\operatorname{disc}(P_7)=49$, so its field
discriminant is at most $49$. The ordinary Minkowski bound
for a totally real cubic field is at most
$(3!/3^3)\sqrt{49}=14/9<2$.
Each ideal class therefore contains an integral ideal of norm
one, proving $\operatorname{Cl}(R)=1$.
This is the elementary application of the retained Minkowski
bound, not a use of a candidate class number.
Altogether every cover used in Steps 2--3 satisfies
\[
                     \operatorname{Cl}(M)[5]=0,
              \qquad [M:\mathbb Q]\in\{3,6,12\}.
 \tag{F5.20}
\]
The initial real-$105$ exploratory calculation and the initial
three-field pilot are not inputs to this twenty-cover proof.

\proofstep{5. Each admitted cover has one prime above five.}
In $R$, the Frobenius class of five has order three in
$(\mathbb Z/7\mathbb Z)^\times/\{1,-1\}$.
Thus $R$ has one prime above five, unramified of degree three.
For a quartic $F_A$, the quadratic twist in (F5.11) is
unramified at five, while $\tau_5^{-t}$ for odd $t$ has full
order-four inertia. The cyclic quartic field therefore has
one prime above five, totally ramified of degree four.
Its local compositum with the unramified cubic field has degree
twelve, equal to $[RF_A:\mathbb Q]$.
Hence every quartic cover has one prime with $(e_5,f_5)=(4,3)$.

For an admitted quadratic radicand divisible by five, the
squarefree polynomial $X^2-D$ is Eisenstein at five. Its
local field is ramified quadratic, so the compositum with
$R$ has one prime with $(e_5,f_5)=(2,3)$.
For the four other admitted radicands $D=2,3,7,42$, the
Legendre symbol $(D/5)$ is $-1$.
The quadratic completion is unramified of degree two.
Its compositum with the unramified cubic field has degree six,
giving one prime with $(e_5,f_5)=(1,6)$.
All these local pairs also occur in the exact class-certificate
outputs, but the uniqueness assertion has been proved directly.

Here is the exact boundary of the one-prime argument.
Let $d_5(\delta)\in\{0,1\}$ be given by
$\delta(5)=(-1)^{d_5(\delta)}$.
Choose a primitive twelfth root $\zeta_{12}$ with
$\zeta_{12}^2=\xi_7(3)$ and $\zeta_{12}^3=\tau_5(2)$;
the common sign makes these choices compatible.
On the field of $\rho$, the five-adic decomposition image
is generated by the exponents
\[
 -3e,\qquad 6\bigl(\epsilon(\delta)+d_5(\delta)-e\bigr)+4j
                              \pmod {12}.
 \tag{F5.21}
\]
The first is inertia. The second uses
$5\equiv3^5\pmod7$ and $\eta_K(5)=-1$ for the unramified
Frobenius lift. If $t$ is odd, inertia supplies the full
two-primary part of order four; if $t=2$, it supplies the
full two-primary part of order two. The second generator
supplies precisely the required cubic part in either case.
For $t=0$, use
\[
 \begin{array}{c|rrrr}
 \delta&1&\chi_{-4}&\chi_8&\chi_{-8}\\\hline
 \epsilon(\delta)&0&1&0&1\\
 d_5(\delta)&0&0&1&1.
 \end{array}\tag{F5.22}
\]
The full image is obtained exactly for the fifteen pairs in
the second line of (F5.1). In the other nine cases the residual
decomposition image has index two. This is only a limitation
of the present method, not a nonvanishing theorem for them.

\proofstep{6. Full strict zero over the cover, then the original rational ring.}
For each admitted pair the cover $M$ kills $\rho$, hence
$a^{\rm nc}|_{G_M}$ takes values in $1+\mathfrak m_O$ on the
unchanged original ring. It is truly unramified at every finite
place away from five. The cyclotomic factor is unramified there;
on that inertia the remaining finite $\psi^{\rm nc}$ has
order dividing six and reduction one. A prime-to-five root
of unity congruent to one is one, since
$(z^n-1)/(z-1)\equiv n\not\equiv0\pmod{\mathfrak m_O}$.

Apply the real-field odd-prime lemma proved in G197 Step 4,
(R7.19)--(R7.20), with $p=5$ and this exact $M$.
For clarity, its full mechanism is as follows.
The image of $G_{M_v}$ in the maximal pro-five quotient of
$G_{M,\{v\}}$ is closed. A proper image would be proper in a
finite five-group quotient and then contained in a normal
subgroup of index five. That quotient would give a cyclic
degree-five extension unramified at every finite place and
split at $v$. Real places also split, since five is odd.
Ordinary Hilbert class-field reciprocity contradicts (F5.20).
Thus the local image is the whole pro-five quotient.

A continuous strict cocycle $z:G_M\to D_O(a^{\rm nc})$ has
finite image. Full local zero gives a single witness $x\in D_O$.
Choose a power of the original uniformizer killing both that
witness and all values of $z$. The resulting homomorphism into
the finite semidirect product of the torsion module with its
near-one scalar image has five-power order.
Away-five strict inertia zero makes this homomorphism trivial
on those inertia groups. It therefore factors through the
maximal pro-five quotient above. Every global image has a
local representative, acting identically on the same witness.
The local coboundary formula is then global.
This proves $\mathcal S_5(M,D_O(a^{\rm nc}))=0$ without a
local-invariant divisibility assumption or a coefficient change.

A rational strict class restricts into that cover strict group,
preserving full local zero at five and zero on all other inertia.
Restriction followed by corestriction is multiplication by
$[M:\mathbb Q]\in\{3,6,12\}$, which is invertible on this
$5$-primary group. This proves (F5.3) on the original $O$.

\proofstep{7. The inert quadratic local field and the critical partner.}
The two characters in (F5.2) restrict to the same action on
$G_K$. Their original integral induced lattice splits as
\[
 \operatorname{Ind}_{G_K}^{G_{\mathbb Q}}
           O(\delta\xi_7^j c_5^{-e})
                  \simeq O(a^{\rm nc})\oplus O(a^{\rm cr}).
 \tag{F5.23}
\]
The determinant of the two-coset change of basis is a unit
times two, and two is a unit in this $O$.
Tensoring with $E/O$ and applying Shapiro therefore gives
an exact decomposition with no lattice-index error.

Five is inert in $K$: its reduction modulo three is two.
There is one actual completion $K_v/\mathbb Q_5$, the
unramified quadratic extension. Local Shapiro at five has
one double coset and identifies the full local group over
that completion with the full rational localization of
the induced coefficient. No split-prime pair is substituted.
At three the ramified quadratic inertia quotient is all
of $C_2$, so inertia Shapiro likewise has one double coset,
with intersection the actual $K$-inertia subgroup.
At every other finite prime the unramified quadratic
Mackey components give exactly the inertia conditions at
all places above it, with repeated conjugate conditions
irrelevant to zero. Real-place $H^1$ is zero at five.
Thus the full strict kernel over $K$ is the direct sum
of the two rational strict kernels. Equation (F5.3)
removes only the noncritical summand and proves (F5.4).

The original localization can include $q,d,269$ and the
other host ramification primes. At an extra prime where
the coefficient is unramified, strict inertia zero is
literal zero of the cocycle. Its finite homomorphism
kills the closed normal subgroup generated by that inertia,
so inflation preserves exactly the strict kernel.
This retains the finite-\'etale/Galois-cohomology dictionary
of (FC.2), including the original completion and normalization.

\proofstep{8. Actual labels and one existing aggregate.}
The actual host contains $\mathbb Q(\zeta_{2^{h+2}})$ and
$\mathbb Q(\zeta_7)$, and $h\geq2$ supplies all four
characters in $\mathcal D$.
The specified original $O_5$ contains their values and all
six values of $\xi_7$, since its character order is
$3\cdot2^h$. The coefficient field also contains $\tau_5$
automatically, but no conductor-five finite character is
adjoined to the host label in (F5.5).
Different displayed rational labels have different
coprime two- and seven-parts. Their ratios have conductor
dividing $56$, and hence cannot equal the kernel character
$\eta_K$, ramified at three. Restriction to $G_K$ preserves
their distinction.

There are twenty-four admitted pairs at each weight not
divisible by four, and fifteen at each multiple of four.
Exactly $\lfloor (T-1)/4\rfloor$ weights are in the second
case, proving (F5.6), including $T=1$.
For each admitted actual triple, (F5.4) identifies its
original strict group with its unique critical rational
component on the same $O_5$ and the same $b_5$.
An arbitrary subset is a subset of the one nonnegative
critical aggregate (CP.2). This proves (F5.7) without a
twenty-four or weight-count multiplier.
The corresponding noncritical terms in (CP.3) are zero
individually by (F5.3), proving (F5.8) for every retained $J$.
The full mixed remainder and the unsupported noncritical
rows remain. This proof neither changes the actual common
window nor supplies the remaining numerical estimate.
\end{proof}

\paragraph{Ownership and source scope.}
Root proposed the route, derived the cover tables and supplied
the original twenty successful class-quotient executions.
Trace independently checked the field/character/local formulas,
wrote this complete original-coefficient proof and performed
the separately recorded exact polynomial and class replays.
The applied PARI flag-one manual, ordinary Hilbert class-field
statement, maximal-$S$ extension construction and Minkowski
bound have their precise source scopes in the manifest.
The local generation argument is fully supplied above; no
new proof of global reciprocity or the certification algorithm
is claimed. The original Artin--Tate/finite-\'etale dictionary
retains SGA1, Expos\'e V, sections 4 and 7, at its preceding
FC scope. Neither an unused real-$105$ candidate nor a complete
unit-basis assertion enters the theorem.
